\section{Independence and rank of molecular graphs}
\label{sec:jacobs}

With the molecule rank formula
$r(G^2) = 3|V| - 6 - \operatorname{def}(\tilde G)$
(\cref{cor:molecule-rank-formula}) in place, this chapter derives two
further results of Jackson--Jord\'an~\cite{jacksonJordan2008} on squares
of graphs. The first is a conjecture of Jacobs~\cite{jacobs1998} on the
independence of molecular graphs, stated as Conjecture~5.1
of~\cite{jacksonJordan2008} and derived there from the rank formula
(their Theorem~5.4); with the rank formula a theorem, it holds
unconditionally: \emph{the square $G^2$ of a graph is independent in the
three-dimensional generic rigidity matroid if and only if $G^2$
satisfies the three-dimensional Laman counting condition}
(\cref{thm:jacobs}). The second is the degree-one rank formula
(Lemma~4.2 of~\cite{jacksonJordan2008}; the tree case is due to
Franzblau~\cite{franzblau2000}): an explicit formula for $r(G^2)$ for
connected graphs with vertices of degree one --- where the rank
formula's minimum-degree hypothesis fails --- by reduction to the
maximal subgraph of minimum degree at least two
(\cref{thm:degree-one-rank}).

The development runs as follows. The three-dimensional Laman condition
and its first consequence --- a graph with Laman square has maximum
degree at most three (Jackson--Jord\'an Lemma~5.2) --- open the chapter.
The substantial combinatorial input is Jackson--Jord\'an's
Theorem~5.3, the counting bound
$|E(G^2)| \le 3|V| - 6 - \operatorname{def}(\tilde G)$ for graphs of
minimum degree at least two with Laman square
(\cref{thm:laman-square-count}); its proof classifies the edges of
$G^2$ against a partition of $V$ attaining the deficiency, using
structural facts about such \emph{tight} partitions proved by
Jackson--Jord\'an in a companion paper (Lemma~3.2
of~\cite{jacksonJordan2008}). Combined with the rank formula, the
counting bound forces $r(G^2) = |E(G^2)|$ for such graphs. The easy
direction of Jacobs' conjecture (independence implies the Laman
condition) is the three-dimensional form of the count bound already
used for Laman's theorem in the plane
(\cref{lem:isSparse-of-rowIndependent-two}). The passage from minimum
degree at least two to arbitrary graphs, and the whole of the
degree-one rank formula, rest on the classical \emph{$0$-extension}
(vertex-addition) lemma in dimension three
(Whiteley~\cite[Lemma~9.1.3]{whiteley1996}; Lemma~3.3
of~\cite{jacksonJordan2008}): attaching a new vertex by at most three
edges to distinct vertices adds exactly that many to the generic rank.

Throughout, $V$ is a finite vertex set, $G$ a simple graph on $V$,
$r = r_3$ the rank function of the three-dimensional generic rigidity
matroid $\mathcal{R}_3(V)$ (\cref{def:genericRank}), and
$\operatorname{def}(\tilde G)$ the $D$-deficiency of
\cref{def:D-deficiency} at $D = 6$, i.e.\
$\operatorname{def}(\tilde G) = \max_P \bigl(6(|P|-1) - 5\,d_G(P)\bigr)$
over partitions $P$ of $V$. An edge set is \emph{independent} when it is
independent in $\mathcal{R}_3(V)$; by
\cref{lem:genericRigidityMatroid-indep-iff} this is row independence at
some (equivalently, any generic) placement.

\subsection{The three-dimensional Laman condition}
\label{sec:jacobs-laman3}

\begin{definition}[The three-dimensional Laman condition]
  \label{def:isLaman3}
  \lean{SimpleGraph.IsLaman3}
  \leanok
  \uses{def:edgesIn}
  A graph $H$ on $V$ is \emph{Laman} (in the three-dimensional counting
  sense of \cite[\S5]{jacksonJordan2008}) if for every $X \subseteq V$
  with $|X| \ge 3$,
  \[
    \bigl|\edgesIn[H]{X}\bigr| \;\le\; 3|X| - 6 .
  \]
\end{definition}

\begin{fmlnote}
  \label{fmlnote:isLaman3-guard}
  The condition is not an instance of the $(k,\ell)$-sparsity family of
  \cref{def:isSparse} at $(k, \ell) = (3, 6)$. That predicate quantifies
  over all $X$ with $\ell \le k|X|$, which at $(3,6)$ includes
  $|X| = 2$, where the bound $3|X| - 6 = 0$ fails on every graph with an
  edge --- while the Laman condition of \cite[\S5]{jacksonJordan2008}
  constrains only $|X| \ge 3$ and holds, for instance, vacuously on
  $K_2$. (At $(2,3)$ the two guards agree, which is why the planar
  \cref{def:isLaman} could be an instance of the family.) The condition
  is therefore formalized as a standalone predicate with the $|X| \ge 3$
  guard; the $(k,\ell)$-sparsity API is unchanged.
\end{fmlnote}

\begin{lemma}[Monotonicity]
  \label{lem:isLaman3-mono}
  \lean{SimpleGraph.IsLaman3.mono_left}
  \leanok
  \uses{def:isLaman3}
  A subgraph of a Laman graph is Laman.
\end{lemma}
\begin{proof}
  \leanok
  \uses{lem:edgesIn-mono}
  The edge counts $\bigl|\edgesIn[H]{X}\bigr|$ only decrease.
\end{proof}

\begin{lemma}[Edge count of a clique]
  \label{lem:clique-edgesIn-count}
  \lean{SimpleGraph.IsClique.ncard_edgesIn}
  \leanok
  \uses{def:edgesIn}
  If $X \subseteq V$ is a clique of a graph $H$ (every two distinct
  vertices of $X$ adjacent), then
  $\bigl|\edgesIn[H]{X}\bigr| = \binom{|X|}{2}$.
\end{lemma}
\begin{proof}
  \leanok
  The edges inside $X$ are exactly the two-element subsets of $X$.
\end{proof}

\begin{lemma}[Maximum degree three; JJ Lemma 5.2]
  \label{lem:isLaman3-degree-le-three}
  \lean{SimpleGraph.IsLaman3.degree_le_three}
  \leanok
  \uses{def:isLaman3, def:square-graph}
  Let $G$ be a graph such that $G^2$ is Laman. Then every vertex of $G$
  has degree at most three.
\end{lemma}
\begin{proof}
  \leanok
  \uses{lem:square-cliques, lem:clique-edgesIn-count}
  For a vertex $v$, the closed neighborhood $N_G[v]$ is a clique of
  $G^2$ (\cref{lem:square-cliques}) on $d_G(v) + 1$ vertices. A clique
  on five vertices spans $\binom{5}{2} = 10 > 3\cdot 5 - 6$ edges
  (\cref{lem:clique-edgesIn-count}), violating the Laman condition, so
  $d_G(v) + 1 \le 4$.
\end{proof}

\subsection{Tight partitions}
\label{sec:jacobs-tight-partitions}

The counting bound is proved against a partition of $V$ attaining the
deficiency. This subsection collects the structural facts about such
partitions that the count uses; they are due to Jackson--Jord\'an
(Lemma~3.2 of~\cite{jacksonJordan2008}, proved in their companion paper
on partitions into maximal subgraphs of zero deficiency). The statements
are given for the $D$-deficiency of \cref{def:D-deficiency} at a general
$D$, matching the formalization; the chapter consumes them at $D = 6$.

\begin{lemma}[Tight partitions exist]
  \label{lem:exists-tight-partition}
  \lean{Graph.IsTightPartition, Graph.exists_isTightPartition}
  \leanok
  \uses{def:D-deficiency}
  Let $G$ be a multigraph on a finite nonempty vertex set. Some
  partition $P$ of $V(G)$ attains the deficiency:
  $\operatorname{def}_{\tilde G}(P) = \operatorname{def}(\tilde G)$.
  Such a partition is called \emph{tight}.
\end{lemma}
\begin{proof}
  \leanok
  The maximum of finitely many values is attained.
\end{proof}

\begin{lemma}[Deficiency under coarsening]
  \label{lem:partitionDef-merge}
  \lean{Graph.crossingEdgesWithin, Graph.partitionDef_merge}
  \leanok
  \uses{def:D-deficiency}
  Let $P$ be a partition of $V(G)$, let $Q \subseteq P$ with
  $|Q| \ge 2$, and let $P'$ be the partition obtained from $P$ by
  merging the members of $Q$ into a single part. Write $e_G(Q)$ for the
  number of edges of $G$ joining distinct members of $Q$. Then
  \[
    \operatorname{def}_{\tilde G}(P') \;=\;
    \operatorname{def}_{\tilde G}(P)
      \;-\; \bigl( D\,(|Q| - 1) - (D-1)\,e_G(Q) \bigr).
  \]
\end{lemma}
\begin{proof}
  \leanok
  Merging removes $|Q| - 1$ parts, and the edges that stop crossing are
  exactly those joining distinct members of $Q$. Read right to left,
  the identity also computes the deficiency change under refining one
  part into several.
\end{proof}

\begin{lemma}[Subfamilies of a tight partition; JJ Lemma 3.2(a)]
  \label{lem:tight-partition-subfamily}
  \lean{Graph.IsTightPartition.subfamily_le}
  \leanok
  \uses{def:D-deficiency, lem:exists-tight-partition}
  Let $P$ be a tight partition of $V(G)$ and $Q \subseteq P$ with
  $|Q| \ge 2$. Then
  \[
    (D-1)\,e_G(Q) \;\le\; D\,(|Q| - 1).
  \]
\end{lemma}
\begin{proof}
  \leanok
  \uses{lem:partitionDef-merge}
  Merging $Q$ (\cref{lem:partitionDef-merge}) changes the deficiency by
  $(D-1)\,e_G(Q) - D(|Q|-1)$; tightness of $P$ makes every change
  nonpositive.
\end{proof}

\begin{fmlnote}
  \label{fmlnote:tight-partition-consumed-forms}
  Jackson--Jord\'an state Lemma~3.2(a) as the nonnegativity
  $\operatorname{def}_H(Q) \ge 0$ of the subfamily's deficiency inside
  the induced subgraph $H = G[\bigcup Q]$, and Lemma~3.2(b) as ``each
  part induces a subgraph of zero deficiency''. The chapter formalizes
  only the forms actually consumed downstream:
  \cref{lem:tight-partition-subfamily} is 3.2(a) rewritten as an
  inequality about $G$ itself (no induced-subgraph object is needed),
  and \cref{lem:tight-partition-parts} is the part-structure dichotomy
  that 3.2(b) is used for, proved directly by splitting a single vertex
  off its part rather than through the zero-deficiency statement.
\end{fmlnote}

\begin{lemma}[Parts of a tight partition; JJ Lemma 3.2(b), consumed form]
  \label{lem:tight-partition-parts}
  \lean{Graph.IsTightPartition.parts}
  \leanok
  \uses{def:D-deficiency}
  Let $G$ be a simple graph, $D \ge 2$, and $P$ a tight partition of
  $V$. Then every part of $P$ is either a singleton or has at least
  three vertices, and in the latter case every vertex of the part has
  at least two neighbors inside its own part.
\end{lemma}
\begin{proof}
  \leanok
  \uses{lem:partitionDef-merge}
  Let $A \in P$ with $|A| \ge 2$ and $v \in A$. Refining $P$ by
  splitting $v$ off as a singleton changes the deficiency by
  $D - (D-1)\,d$, where $d$ is the number of neighbors of $v$ inside
  $A$ (\cref{lem:partitionDef-merge}, read as a refinement); tightness
  forces $D \le (D-1)\,d$, hence $d \ge 2$. A two-vertex part is then
  impossible: each of its vertices would need two neighbors inside the
  part, but a simple graph has at most one edge on two vertices.
\end{proof}

\begin{lemma}[Edge multiplicity between parts of a tight partition]
  \label{lem:tight-partition-cross-pair-mult}
  \lean{Graph.IsTightPartition.crossingEdgesWithin_pair_le_one}
  \leanok
  \uses{def:D-deficiency}
  Let $G$ be a simple graph and $P$ a tight partition of $V$ with
  $D \ge 3$. Then any two distinct parts of $P$ are joined by at most
  one edge of $G$.
\end{lemma}
\begin{proof}
  \leanok
  \uses{lem:tight-partition-subfamily}
  An instance of \cref{lem:tight-partition-subfamily}: two edges
  between parts $A \ne B$ give $e_G(Q) \ge 2$ for $Q = \{A, B\}$,
  contradicting $(D-1)\cdot 2 \le D$ for $D \ge 3$.
\end{proof}

\begin{lemma}[Uniqueness of the common neighbor of a cross pair]
  \label{lem:tight-partition-cross-pair-nbr}
  \lean{Graph.IsTightPartition.eq_of_common_nbr}
  \leanok
  \uses{def:D-deficiency}
  Let $G$ be a simple graph and $P$ a tight partition of $V$ with
  $D \ge 5$. Then any two nonadjacent vertices lying in distinct parts
  have at most one common neighbor.
\end{lemma}
\begin{proof}
  \leanok
  \uses{lem:tight-partition-subfamily, lem:tight-partition-cross-pair-mult}
  Let $u, w$ be nonadjacent in distinct parts with common
  neighbors $v \ne v'$; the four edges $uv, vw, uv', v'w$ join the at
  most four parts containing $u, v, w, v'$, and a case analysis on
  which of these parts coincide produces, in every case, a subfamily
  $Q$ with either two edges between two of its members (excluded by
  \cref{lem:tight-partition-cross-pair-mult}) or $e_G(Q) \ge |Q|$ with
  $|Q| \le 4$, contradicting \cref{lem:tight-partition-subfamily} for
  $D \ge 5$.
\end{proof}

\subsection{The counting bound}
\label{sec:jacobs-counting}

Fix now a simple graph $G$ of minimum degree at least two whose square
is Laman, and a tight partition $P$ of $V$ at $D = 6$. By
\cref{lem:isLaman3-degree-le-three} every vertex has degree two or
three, and by \cref{lem:tight-partition-parts} every part of $P$ is a
singleton or has at least three vertices with minimum in-part degree
two. Following \cite[\S5]{jacksonJordan2008}, call an edge
$uw \in E(G^2)$ joining distinct parts a \emph{cross edge} if
$uw \notin E(G)$; its endpoints have a unique common neighbor $v$
(\cref{lem:tight-partition-cross-pair-nbr}), and the cross edge is
\emph{normal} if $v$ lies in the same part as $u$ or as $w$, and
\emph{special} otherwise.

\begin{lemma}[Classification of the edges of the square]
  \label{lem:square-cross-classification}
  \lean{SimpleGraph.square_edgeSet_eq_union, SimpleGraph.disjoint_inPart_gCross,
    SimpleGraph.disjoint_inPart_cross, SimpleGraph.disjoint_gCross_cross,
    SimpleGraph.squareCrossEdges_eq_union,
    SimpleGraph.squareNormalCrossEdges_disjoint_special,
    SimpleGraph.squareGCrossEdges_ncard_eq_crossingEdges,
    SimpleGraph.squareSpecialCrossEdges_singleton_part,
    SimpleGraph.squareNormalCrossEdges_part_three_le}
  \leanok
  \uses{def:square-graph, def:D-deficiency}
  Let $G$, $P$ be as above. Then $E(G^2)$ is the disjoint union of four
  classes: the edges inside a part; the edges of $G$ joining distinct
  parts (numbering $d_G(P)$); the normal cross edges; and the special
  cross edges. Moreover the common neighbor of a special cross edge
  forms a singleton part, and the common neighbor of a normal cross
  edge lies in a part of at least three vertices, together with exactly
  one endpoint of the edge.
\end{lemma}
\begin{proof}
  \leanok
  \uses{lem:tight-partition-cross-pair-nbr, lem:isLaman3-degree-le-three,
    lem:tight-partition-parts}
  An edge of $G^2$ either lies inside a part or joins distinct parts,
  and in the latter case is an edge of $G$ or a cross edge; the
  normal/special dichotomy is well defined because the common neighbor
  is unique (\cref{lem:tight-partition-cross-pair-nbr}). For a special
  cross edge $uw$ with common neighbor $v$: both $u$ and $w$ lie
  outside $v$'s part, so $v$ has two neighbors outside its part; if
  that part were not a singleton, $v$ would also have two neighbors
  inside it (\cref{lem:tight-partition-parts}), exceeding the degree
  bound three (\cref{lem:isLaman3-degree-le-three}). For a normal
  cross edge, the common neighbor shares a part with exactly one
  endpoint (the endpoints lie in distinct parts), and a part
  containing two adjacent vertices is not a singleton, hence has at
  least three vertices (\cref{lem:tight-partition-parts}).
\end{proof}

\begin{lemma}[Special cross edges at a singleton part; JJ eq.\ (5), (7)]
  \label{lem:singleton-part-neighborhood}
  \lean{SimpleGraph.isClique_neighborSet_square,
    SimpleGraph.ncard_edgesIn_neighborSet_square,
    SimpleGraph.IsSquareTightPartition.mem_squareSpecialCrossEdges_of_singleton_part}
  \leanok
  \uses{def:square-graph, def:D-deficiency}
  Let $G$, $P$ be as above and let $\{v\} \in P$ be a singleton part.
  The neighborhood $N_G(v)$ is a clique of $G^2$ all of whose edges are
  special cross edges with common neighbor $v$, and
  \[
    \bigl|\edgesIn[G^2]{N_G(v)}\bigr| \;=\; 2\,d_G(v) - 3 .
  \]
\end{lemma}
\begin{proof}
  \leanok
  \uses{lem:square-cliques, lem:clique-edgesIn-count,
    lem:isLaman3-degree-le-three, lem:tight-partition-cross-pair-mult,
    lem:tight-partition-cross-pair-nbr, lem:tight-partition-subfamily}
  $N_G(v)$ is a clique of $G^2$ (\cref{lem:square-cliques}). Two
  neighbors $u, w$ of $v$ lie in parts distinct from $\{v\}$ and from
  each other --- two neighbors in one part $A$ would put two edges
  between $A$ and $\{v\}$ (\cref{lem:tight-partition-cross-pair-mult}) ---
  and are nonadjacent in $G$: a triangle $uvw$ would give a three-part
  subfamily with three edges between its members, excluded by
  \cref{lem:tight-partition-subfamily} for $D \ge 4$. So each pair is a
  cross edge with common neighbor $v$ (unique
  by \cref{lem:tight-partition-cross-pair-nbr}), special because $v$'s
  part $\{v\}$ contains neither endpoint. Since
  $d_G(v) \in \{2, 3\}$ (\cref{lem:isLaman3-degree-le-three}, minimum
  degree two), the clique count $\binom{d_G(v)}{2}$
  (\cref{lem:clique-edgesIn-count}) equals $2\,d_G(v) - 3$ in both
  cases.
\end{proof}

\begin{lemma}[Special cross edges arise from a unique singleton part]
  \label{lem:singleton-part-converse}
  \lean{SimpleGraph.exists_unique_singleton_part_of_mem_squareSpecialCrossEdges}
  \leanok
  \uses{def:square-graph, def:D-deficiency}
  Let $G$, $P$ be as above. Every special cross edge arises, as in
  \cref{lem:singleton-part-neighborhood}, from exactly one singleton
  part.
\end{lemma}
\begin{proof}
  \leanok
  \uses{lem:singleton-part-neighborhood, lem:square-cross-classification,
    lem:tight-partition-cross-pair-nbr}
  A special cross edge's common neighbor is a singleton part
  (\cref{lem:square-cross-classification}), and it is unique
  (\cref{lem:tight-partition-cross-pair-nbr}), so the assignment of
  \cref{lem:singleton-part-neighborhood} is a partition of the special
  cross edges.
\end{proof}

\begin{lemma}[The root of a normal cross edge; JJ eq.\ (6) setup]
  \label{lem:normal-cross-count-rooted}
  \lean{SimpleGraph.squareNormalCrossEdgesRootedAt,
    SimpleGraph.IsSquareTightPartition.rootedAt_inj}
  \leanok
  \uses{def:square-graph, def:D-deficiency}
  Let $G$, $P$ be as above. Say a normal cross edge is \emph{rooted} at
  the part containing its common neighbor. Every normal cross edge is
  rooted at exactly one part, and that part has at least three vertices.
\end{lemma}
\begin{proof}
  \leanok
  \uses{lem:square-cross-classification, lem:tight-partition-cross-pair-nbr}
  The common neighbor is unique
  (\cref{lem:tight-partition-cross-pair-nbr}), so the root is well
  defined; and by the classification (\cref{lem:square-cross-classification})
  the common neighbor of a normal cross edge lies in a part of at least
  three vertices, together with exactly one endpoint of the edge.
\end{proof}

\begin{lemma}[The local count; JJ eq.\ (6), local half]
  \label{lem:normal-cross-count-local}
  \lean{SimpleGraph.IsSquareTightPartition.ncard_inPartNeighbors_eq_two}
  \leanok
  \uses{def:square-graph, def:D-deficiency}
  Let $G$, $P$ be as above, let $A$ be a part with $|A| \ge 3$, and let
  $v \in A$ lie on a crossing edge $vw$ of $G$ (an edge with
  $w \notin A$). Then $v$ has exactly two neighbors inside $A$.
\end{lemma}
\begin{proof}
  \leanok
  \uses{lem:isLaman3-degree-le-three, lem:tight-partition-parts}
  The in-part degree of $v$ is at least two
  (\cref{lem:tight-partition-parts}) and its total degree at most three
  (\cref{lem:isLaman3-degree-le-three}); the neighbors of $v$ split into
  those inside $A$ and those outside, and $w$ witnesses that the outside
  set is nonempty, so exactly two neighbors of $v$ lie inside $A$.
\end{proof}

\begin{lemma}[Producing a normal cross edge; JJ eq.\ (6), construction]
  \label{lem:normal-cross-count-producer}
  \lean{SimpleGraph.IsSquareTightPartition.mk_mem_squareNormalCrossEdgesRootedAt}
  \leanok
  \uses{def:square-graph, def:D-deficiency}
  Let $G$, $P$ be as above, and let $v$ have an in-part neighbor $u$
  ($f u = f v$) and an out-of-part neighbor $w$ ($f w \ne f v$). Then
  $uw$ is a normal cross edge rooted at $v$'s part, and $v$ is its
  unique common neighbor.
\end{lemma}
\begin{proof}
  \leanok
  \uses{lem:square-cliques, lem:tight-partition-cross-pair-mult,
    lem:tight-partition-cross-pair-nbr}
  Both $u, w \in N_G(v)$, so $uw \in E(G^2)$
  (\cref{lem:square-cliques}); it joins distinct parts
  ($f u = f v \ne f w$) and is not an edge of $G$ --- otherwise $vw$ and
  $uw$ would be two edges between $v$'s part and $w$'s
  (\cref{lem:tight-partition-cross-pair-mult}). Its common neighbor $v$
  lies in $u$'s part, so the edge is normal and rooted there, and the
  common neighbor is unique (\cref{lem:tight-partition-cross-pair-nbr}).
\end{proof}

\begin{lemma}[Two rooted normal cross edges per crossing edge; JJ eq.\ (6), fiber]
  \label{lem:normal-cross-count-fiber}
  \lean{SimpleGraph.IsSquareTightPartition.ncard_normalCrossEdges_of_crossing_eq_two}
  \leanok
  \uses{def:square-graph, def:D-deficiency}
  Let $G$, $P$ be as above, let $A$ be a part with $|A| \ge 3$, and let
  $vw$ be a crossing edge of $G$ with $v \in A$, $w \notin A$. Then the
  normal cross edges $uw$ produced from $vw$ by the in-part neighbors
  $u \in N_G(v) \cap A$ (\cref{lem:normal-cross-count-producer}) number
  exactly two.
\end{lemma}
\begin{proof}
  \leanok
  \uses{lem:normal-cross-count-producer, lem:normal-cross-count-local}
  The map $u \mapsto uw$ is injective on $N_G(v) \cap A$ (the endpoint
  $w$ is shared and $u \ne w$, as $u \in A \not\ni w$), so its image has
  the same size as $N_G(v) \cap A$, which is two
  (\cref{lem:normal-cross-count-local}); each image edge is a rooted
  normal cross edge by \cref{lem:normal-cross-count-producer}.
\end{proof}

\begin{definition}[Edges cut by a part]
  \label{def:g-cut-edges}
  \lean{SimpleGraph.gCutEdges}
  \leanok
  \uses{def:D-deficiency}
  For a partition $P$ of $V$ (a labeling $f$) and a part $A$, write
  $d_G(A)$ for the number of edges of $G$ with \emph{exactly one}
  endpoint in $A$: edges $xy \in E(G)$ with $x \in A$ and $y \notin A$.
\end{definition}

\begin{lemma}[The rooted normal cross edges partition over the crossing edges; JJ eq.\ (6)]
  \label{lem:normal-cross-count-partition}
  \lean{SimpleGraph.squareCutPairs, SimpleGraph.ncard_squareCutPairs_eq_gCutEdges,
    SimpleGraph.IsSquareTightPartition.squareNormalCrossEdgesRootedAt_eq_biUnion,
    SimpleGraph.IsSquareTightPartition.squareCutPairs_pairwiseDisjoint}
  \leanok
  \uses{def:square-graph, def:D-deficiency, def:g-cut-edges}
  Let $G$, $P$ be as above and $A$ a part. Index the \emph{oriented}
  crossing edges of $A$ by the pairs $(v, w)$ with $v \in A$,
  $w \notin A$, $vw \in E(G)$; the map $(v, w) \mapsto vw$ is a bijection
  onto the $d_G(A)$ cut edges of $A$ (\cref{def:g-cut-edges}). The normal
  cross edges rooted at $A$ are the union, over these oriented crossing
  edges, of the fibers $\{uw : u \in N_G(v) \cap A\}$, and fibers over
  distinct oriented crossing edges are disjoint.
\end{lemma}
\begin{proof}
  \leanok
  \uses{lem:normal-cross-count-producer, lem:square-cross-classification,
    lem:tight-partition-cross-pair-nbr}
  The bijection with cut edges holds because a cut edge has a unique
  endpoint in $A$, fixing the orientation. Each fiber edge $uw$ is a
  rooted normal cross edge (\cref{lem:normal-cross-count-producer});
  conversely a rooted normal cross edge has a common neighbor $v \in A$
  and, by the classification (\cref{lem:square-cross-classification}),
  exactly one endpoint in $A$, so it lies in the fiber of
  $(v, \text{other endpoint})$. Disjointness: an edge in two fibers has
  the same common neighbor (\cref{lem:tight-partition-cross-pair-nbr})
  and the same out-of-part endpoint, so the oriented crossing edges
  coincide.
\end{proof}

\begin{lemma}[Normal cross edges rooted at a part; JJ eq.\ (6)]
  \label{lem:normal-cross-count}
  \lean{SimpleGraph.IsSquareTightPartition.ncard_normalCrossEdgesRootedAt_eq_two_mul_gCutEdges}
  \leanok
  \uses{def:square-graph, def:D-deficiency, def:g-cut-edges}
  Let $G$, $P$ be as above. For each part $A$ with $|A| \ge 3$ the
  normal cross edges rooted at $A$ (\cref{lem:normal-cross-count-rooted})
  number exactly $2\,d_G(A)$ (\cref{def:g-cut-edges}).
\end{lemma}
\begin{proof}
  \leanok
  \uses{lem:normal-cross-count-partition, lem:normal-cross-count-fiber}
  By \cref{lem:normal-cross-count-partition} the rooted normal cross
  edges are a disjoint union of fibers, one per oriented crossing edge,
  and there are $d_G(A)$ oriented crossing edges; each fiber has exactly
  two elements (\cref{lem:normal-cross-count-fiber}). Summing the
  constant $2$ over the $d_G(A)$ fibers gives $2\,d_G(A)$.
\end{proof}

\begin{fmlnote}
  \label{fmlnote:normal-cross-split}
  As anticipated, this node decomposed during formalization. Four
  pieces are separated out: the well-definedness of the root
  (\cref{lem:normal-cross-count-rooted}); the local count --- each
  crossing edge at $A$ contributes exactly two pairs, i.e.\ its endpoint
  in $A$ has exactly two in-part neighbors
  (\cref{lem:normal-cross-count-local}); the construction of a rooted
  normal cross edge from such a pair
  (\cref{lem:normal-cross-count-producer}), the map $(vw, u) \mapsto uw$;
  and the disjoint partition of the rooted normal cross edges into the
  fibers over $A$'s $d_G(A)$ crossing edges
  (\cref{lem:normal-cross-count-partition}). The residual arithmetic in
  \cref{lem:normal-cross-count} --- the disjoint-union count, summing
  the constant fiber size $2$ over the $d_G(A)$ fibers --- closes the
  node directly.
\end{fmlnote}

\begin{theorem}[The counting bound for Laman squares; JJ Theorem 5.3]
  \label{thm:laman-square-count}
  \lean{SimpleGraph.laman_square_count}
  \leanok
  \uses{def:isLaman3, def:square-graph, def:D-deficiency}
  Let $G$ be a simple graph of minimum degree at least two on a finite
  nonempty vertex set $V$, and suppose $G^2$ is Laman. Then
  \[
    |E(G^2)| \;\le\; 3|V| - 6 - \operatorname{def}(\tilde G).
  \]
\end{theorem}
\begin{proof}
  \leanok
  \uses{lem:exists-tight-partition, lem:square-cross-classification,
    lem:singleton-part-neighborhood, lem:singleton-part-converse,
    lem:normal-cross-count, lem:tight-partition-parts}
  Let $P = \{P_1, \dots, P_t\}$ be a tight partition
  (\cref{lem:exists-tight-partition}) --- the argument below is uniform
  in $P$, with no separate case for $\operatorname{def}(\tilde G) = 0$.
  Summing the classification of \cref{lem:square-cross-classification}
  with the counts of \cref{lem:singleton-part-neighborhood} and
  \cref{lem:normal-cross-count}, and bounding the edges of $G^2$ inside
  a part $P_i$ with $|P_i| \ge 3$ by the Laman condition (a singleton
  part spans none):
  \[
    |E(G^2)| \;\le\; \sum_{i=1}^{t} \bigl(3|P_i| - 6\bigr)
      + d_G(P) + \sum_{i=1}^{t} 2\,d_G(P_i),
  \]
  where for a singleton part the two summands
  $(3\cdot 1 - 6) + 2\,d_G(P_i) = 2\,d_G(P_i) - 3$ reproduce the
  special-edge count of \cref{lem:singleton-part-neighborhood} (using
  $d_G(P_i) = \deg_G(v)$ for the singleton part $P_i = \{v\}$, since
  every edge at $v$ crosses). Since $\sum_i |P_i| = |V|$ and
  $\sum_i d_G(P_i) = 2\,d_G(P)$ (each crossing edge is counted from
  both sides), the right-hand side is
  $3|V| - 6t + 5\,d_G(P) = 3|V| - 6 - \operatorname{def}_{\tilde G}(P)
  = 3|V| - 6 - \operatorname{def}(\tilde G)$.
\end{proof}

\subsection{Independence implies the Laman condition}
\label{sec:jacobs-easy}

The easy direction of Jacobs' conjecture holds for every graph on $V$,
square or not, by the same argument as in the plane
(\cref{lem:isSparse-of-rowIndependent-two}): a row-independent edge set
cannot exceed the rank of any framework containing it.

\begin{lemma}[The Laman condition from row independence, dimension three]
  \label{lem:isLaman3-of-rowIndependent}
  \lean{SimpleGraph.isLaman3_of_edgeSetRowIndependent_dim_three}
  \leanok
  \uses{def:isLaman3, def:edgeSet-rowIndependent}
  Let $p : V \to \R^3$ be a placement whose restriction to every subset
  of at least four vertices affinely spans $\R^3$, and let $H$ be a
  graph on $V$ whose edge set is row-independent at $p$. Then $H$ is
  Laman.
\end{lemma}
\begin{proof}
  \leanok
  \uses{lem:rigidityMap-finrank-range-le-of-affinelySpanning}
  Fix $X \subseteq V$ with $|X| \ge 3$. For $|X| = 3$ the claim is the
  trivial count $\binom{3}{2} = 3 = 3\cdot 3 - 6$. For $|X| \ge 4$,
  transport the rows of the edges inside $X$ to the induced subgraph
  $H[X]$ at the restricted placement, as in
  \cref{lem:isSparse-of-rowIndependent-two}; row independence gives
  $\bigl|\edgesIn[H]{X}\bigr| =
  \operatorname{rank} R(H[X], p|_X)$, and since $p|_X$ affinely spans,
  the kernel of the rigidity map contains the six-dimensional space of
  infinitesimal rigid motions
  (\cref{lem:rigidityMap-finrank-range-le-of-affinelySpanning}), so the
  rank is at most $3|X| - 6$.
\end{proof}

\begin{corollary}[Independent edge sets are Laman]
  \label{cor:genericMatroid-indep-isLaman3}
  \lean{SimpleGraph.isLaman3_of_genericRigidityMatroid_indep}
  \leanok
  \uses{def:isLaman3, def:genericRigidityMatroid}
  Let $H$ be a graph on $V$ with $E(H)$ independent in
  $\mathcal{R}_3(V)$. Then $H$ is Laman.
\end{corollary}
\begin{proof}
  \leanok
  \uses{lem:genericRigidityMatroid-indep-iff,
    lem:isLaman3-of-rowIndependent, lem:exists-generic-general-position,
    def:general-position-placement}
  Choose a placement simultaneously generic for row independence and in
  general position up to order four
  (\cref{lem:exists-generic-general-position}). Independence makes
  $E(H)$ row-independent at it
  (\cref{lem:genericRigidityMatroid-indep-iff}), and general position
  makes every subset of at least four points affinely spanning
  (\cref{def:general-position-placement}: four of them are affinely
  independent), so \cref{lem:isLaman3-of-rowIndependent} applies.
\end{proof}

\subsection{Zero-extensions in dimension three}
\label{sec:jacobs-zero-extension}

The reduction of Jacobs' conjecture to the minimum-degree-two case, and
the degree-one rank formula, both run on the classical $0$-extension
lemma. Whiteley's form (\cite[Lemma~9.1.3]{whiteley1996}) is a
fixed-placement statement: attaching a new vertex on three edges,
placed non-coplanar with its three neighbors, preserves row
independence and adds exactly three to the rank; Jackson--Jord\'an
(Lemma~3.3 of~\cite{jacksonJordan2008}) cite the independence statement
for a new vertex joined to $s \le 3$ distinct old vertices. On a fixed
vertex type the operation compares a graph $H$ with the graph
$H - E_H(v)$ obtained by deleting the edges at a vertex $v$. Two
corollaries below package the generic consequences. At $d_H(v) \le 3$
the rank grows by exactly $d_H(v)$, and independence is preserved in
both directions, with no further hypothesis. At arbitrary degree the
rank addition $\min(3, d_H(v))$ requires the neighborhood $N_H(v)$ to
be a clique of $H$. The hypothesis is not removable: $K_{1,4}$ has rank
four, not three, since each of its four rows has a coordinate block, at
its own leaf, met by no other row. It is also the hypothesis
Jackson--Jord\'an invoke where they apply the lemma at unbounded
degree: ``the neighbour sets \dots\ induce complete (and hence rigid)
subgraphs'' (proof of Theorem~4.3 of~\cite{jacksonJordan2008}). Both
applications below have it, since the square-neighborhood of a
degree-one vertex lies in a closed neighborhood of $G$, a clique of
$G^2$ (\cref{lem:square-leaf-neighborhood}).

\begin{lemma}[Row-independence lift for a vertex of degree at most three]
  \label{lem:zero-extension-rowIndependent}
  \lean{SimpleGraph.zero_extension_edgeSetRowIndependent_lift}
  \leanok
  \uses{def:edgeSet-rowIndependent}
  Let $H$ be a graph on $V$, let $v$ be a vertex of $H$ of degree
  $s \le 3$ with neighbors $u_1, \dots, u_s$, and let
  $H' = H - E_H(v)$. Suppose $E(H')$ is row-independent at a placement
  $p' : V \to \R^3$ whose images $p'(u_1), \dots, p'(u_s)$ are affinely
  independent. Then there is a placement $p$, agreeing with $p'$ away
  from $v$, at which $E(H)$ is row-independent.
\end{lemma}
\begin{proof}
  \leanok
  Move only $v$: the rows of $E(H')$ do not involve $v$'s coordinates
  and stay independent. Choose $p(v)$ outside the affine hull of
  $p'(u_1), \dots, p'(u_s)$, a proper affine subspace by the affine
  independence; then the $s + 1$ points $p(v), p'(u_1), \dots, p'(u_s)$
  are affinely independent, so the difference vectors $p(v) - p'(u_i)$
  are linearly independent. The $s$ new rows are then independent from
  each other and from the old ones, by the test-motion argument of the
  planar lifts (\cref{lem:typeI-rowIndependent-lift},
  \cref{lem:pendant-rowIndependent-lift}): a motion supported at $v$
  annihilates every row of $E(H')$, and evaluating a dependence at such
  motions forces the new coefficients to vanish.
\end{proof}

The affine-independence hypothesis is likewise not removable: $K_5$
minus an edge $u_1u_2$ admits a row-independent placement with
$p'(u_1) = p'(u_2)$, and after adjoining a vertex $v$ joined to $u_1$
and $u_2$, every placement agreeing with $p'$ away from $v$ is
row-dependent. The corollaries below discharge the hypothesis by taking
$p'$ simultaneously generic for row independence and in general
position (\cref{lem:exists-generic-general-position}).

\begin{corollary}[$0$-extension at degree at most three]
  \label{cor:zero-extension-degree-le-three}
  \uses{def:genericRank, def:genericRigidityMatroid}
  Let $H$ be a graph on $V$ and $v$ a vertex of $H$ with
  $d_H(v) \le 3$. Then
  \[
    r(H) \;=\; r\bigl(H - E_H(v)\bigr) + d_H(v),
  \]
  and $E(H)$ is independent in $\mathcal{R}_3(V)$ if and only if
  $E(H - E_H(v))$ is.
\end{corollary}
\begin{proof}
  \uses{lem:zero-extension-rowIndependent,
    lem:genericRigidityMatroid-indep-iff,
    lem:exists-generic-general-position, def:general-position-placement}
  For the lower bound, a maximal independent subset of
  $E(H - E_H(v))$ is row-independent at a placement simultaneously
  generic for row independence and in general position
  (\cref{lem:exists-generic-general-position}); general position makes
  the images of the $d_H(v) \le 3$ neighbors of $v$ affinely
  independent (\cref{def:general-position-placement}), so
  \cref{lem:zero-extension-rowIndependent}, applied to the subgraph
  carrying that subset and the edges at $v$, adjoins the edges at $v$
  (\cref{lem:genericRigidityMatroid-indep-iff}). For the upper bound,
  adjoining $d_H(v)$ edges raises the rank by at most $d_H(v)$. The
  independence equivalence follows since
  $|E(H)| = |E(H - E_H(v))| + d_H(v)$ and independence of a finite
  edge set is the equality of its rank with its cardinality.
\end{proof}

\begin{corollary}[Rank of a $0$-extension at a clique neighborhood]
  \label{cor:zero-extension-clique-rank}
  \uses{def:genericRank}
  Let $H$ be a graph on $V$ and $v$ a vertex of $H$ whose neighborhood
  $N_H(v)$ is a clique of $H$. Then
  \[
    r(H) \;=\; r\bigl(H - E_H(v)\bigr) + \min\bigl(3,\, d_H(v)\bigr).
  \]
\end{corollary}
\begin{proof}
  \uses{cor:zero-extension-degree-le-three,
    cor:genericMatroid-indep-isLaman3}
  For $d_H(v) \le 3$ this is \cref{cor:zero-extension-degree-le-three}.
  For $d_H(v) \ge 3$, fix three neighbors $u_1, u_2, u_3$ and let
  $H_3$ be $H - E_H(v)$ together with the edges $vu_1, vu_2, vu_3$;
  then $r(H) \ge r(H_3) = r(H - E_H(v)) + 3$ by
  \cref{cor:zero-extension-degree-le-three}. For the reverse
  inequality it suffices that every further edge $vw$,
  $w \in N_H(v) \setminus \{u_1, u_2, u_3\}$, lies in the closure of
  $E(H_3)$. The five vertices $v, u_1, u_2, u_3, w$ span a $K_5$ in
  $H$ whose edges other than $vw$ all lie in $E(H_3)$: the six edges
  among $u_1, u_2, u_3, w$ are edges of the clique $N_H(v)$, none of
  them incident to $v$. That $K_5$ minus $vw$ is independent —
  starting from the empty graph, add its nine edges vertex by vertex,
  each step a $0$-extension of degree at most three
  (\cref{cor:zero-extension-degree-le-three}) — while the full $K_5$
  is dependent, since ten edges on five vertices violate the Laman
  count (\cref{cor:genericMatroid-indep-isLaman3}). Hence $vw$ lies in
  the closure of the other nine edges, and
  $r(H) = r(H_3) = r(H - E_H(v)) + 3$.
\end{proof}

\begin{lemma}[Independence and rank under vertex-set restriction]
  \label{lem:genericMatroid-induce-transport}
  \uses{def:genericRigidityMatroid, def:genericRank}
  Let $d \ge 0$, let $S \subseteq V$ and let $H$ be a graph on $S$.
  Then $E(H)$ is independent in $\mathcal{R}_d(S)$ if and only if its
  image under the inclusion $S \hookrightarrow V$ is independent in
  $\mathcal{R}_d(V)$, and the rank of $E(H)$ in $\mathcal{R}_d(S)$
  equals the rank of its image in $\mathcal{R}_d(V)$.
\end{lemma}
\begin{proof}
  \uses{lem:genericRigidityMatroid-indep-iff}
  Both independences say the edge set is row-independent at some
  placement (\cref{lem:genericRigidityMatroid-indep-iff}). A placement
  of $S$ extends to one of $V$ arbitrarily, and a placement of $V$
  restricts to one of $S$; in either direction the rows of the edges
  inside $S$ are unchanged up to the coordinate inclusion, as in the
  row transport of \cref{lem:isSparse-of-rowIndependent-two}. The rank
  identity follows: the image of a maximal independent subset of
  $E(H)$ is independent of the same cardinality, and every independent
  subset of the image pulls back.
\end{proof}

The reduction also uses three facts about squares and induced
subgraphs.

\begin{lemma}[Deleting the edges at a vertex of degree at most one
  commutes with squaring]
  \label{lem:square-delete-degree-le-one}
  \lean{SimpleGraph.square_deleteIncidenceSet_of_degree_le_one}
  \leanok
  \uses{def:square-graph}
  Let $G$ be a graph and $v$ a vertex of degree at most one in $G$.
  Then
  \[
    \bigl(G - E_G(v)\bigr)^2 \;=\; G^2 - E_{G^2}(v).
  \]
\end{lemma}
\begin{proof}
  \leanok
  A path of length two between vertices $x, y \ne v$ cannot pass
  through $v$, which has at most one neighbor; so the squares agree
  away from $v$, and both sides have no edge at $v$.
\end{proof}

\begin{lemma}[The square-neighborhood of a degree-one vertex]
  \label{lem:square-leaf-neighborhood}
  \lean{SimpleGraph.neighborSet_square_of_degree_eq_one,
    SimpleGraph.ncard_neighborSet_square_of_degree_eq_one,
    SimpleGraph.isClique_neighborSet_square_of_degree_eq_one}
  \leanok
  \uses{def:square-graph}
  Let $v$ be a vertex of degree one in $G$, with neighbor $u$. Then
  \[
    N_{G^2}(v) \;=\; \bigl(N_G(u) \cup \{u\}\bigr) \setminus \{v\},
  \]
  so $d_{G^2}(v) = d_G(u)$, and $N_{G^2}(v)$ is a clique of $G^2$.
\end{lemma}
\begin{proof}
  \leanok
  \uses{lem:square-cliques}
  A vertex $w \ne v$ is a square-neighbor of $v$ exactly when it is
  adjacent to $v$ — so $w = u$ — or shares a common neighbor with
  $v$, which can only be $u$; conversely $u$ and every other neighbor
  of $u$ are within distance two of $v$. The set is contained in the
  closed neighborhood $N_G[u]$, a clique of $G^2$
  (\cref{lem:square-cliques}).
\end{proof}

\begin{lemma}[Squaring commutes with restriction to the support]
  \label{lem:square-induce-support}
  \lean{SimpleGraph.square_induce_of_support_subset, SimpleGraph.square_support_subset}
  \leanok
  \uses{def:square-graph}
  Let $S \subseteq V$ contain every non-isolated vertex of $G$. Then
  the square of the induced subgraph $G[S]$ is the induced subgraph of
  $G^2$ on $S$, and every edge of $G^2$ has both endpoints in $S$.
\end{lemma}
\begin{proof}
  \leanok
  Both endpoints of an edge of $G^2$ have an edge of $G$, so lie in
  $S$; and a path of length at most two between vertices of $S$ has
  its middle vertex non-isolated, hence in $S$, so adjacency in $G^2$
  and in $G[S]^2$ agree on $S$.
\end{proof}

\begin{lemma}[The Laman condition restricts to induced subgraphs]
  \label{lem:isLaman3-induce}
  \lean{SimpleGraph.IsLaman3.induce}
  \leanok
  \uses{def:isLaman3}
  Let $G$ be Laman and $S \subseteq V$. Then the induced subgraph
  $G[S]$ is Laman.
\end{lemma}
\begin{proof}
  \leanok
  A set $X \subseteq S$ of at least three vertices spans in $G[S]$
  exactly the edges it spans in $G$.
\end{proof}

\subsection{Jacobs' conjecture}
\label{sec:jacobs-theorem}

\begin{theorem}[Jacobs' conjecture, minimum degree two]
  \label{thm:jacobs-min-degree-two}
  \lean{SimpleGraph.jacobs_min_degree_two}
  \leanok
  \uses{def:isLaman3, def:square-graph, def:genericRigidityMatroid}
  Let $G$ be a simple graph of minimum degree at least two on a finite
  nonempty vertex set $V$. If $G^2$ is Laman, then $E(G^2)$ is
  independent in $\mathcal{R}_3(V)$.
\end{theorem}
\begin{proof}
  \leanok
  \uses{thm:laman-square-count, cor:molecule-rank-formula}
  By \cref{thm:laman-square-count} and the rank formula
  (\cref{cor:molecule-rank-formula}),
  \[
    |E(G^2)| \;\le\; 3|V| - 6 - \operatorname{def}(\tilde G)
    \;=\; r(G^2) \;\le\; |E(G^2)|,
  \]
  so $r(G^2) = |E(G^2)|$, which is independence.
\end{proof}

\begin{theorem}[Jacobs' conjecture; JJ Conjecture 5.1, Theorem 5.4]
  \label{thm:jacobs}
  \uses{def:isLaman3, def:square-graph, def:genericRigidityMatroid}
  Let $G$ be a simple graph on a finite vertex set $V$. Then $E(G^2)$
  is independent in $\mathcal{R}_3(V)$ if and only if $G^2$ is Laman.
\end{theorem}
\begin{proof}
  \uses{cor:genericMatroid-indep-isLaman3, thm:jacobs-min-degree-two,
    lem:isLaman3-degree-le-three, lem:isLaman3-mono,
    cor:zero-extension-degree-le-three, lem:genericMatroid-induce-transport,
    lem:square-delete-degree-le-one, lem:square-leaf-neighborhood,
    lem:square-induce-support, lem:isLaman3-induce}
  Independence implies the Laman condition by
  \cref{cor:genericMatroid-indep-isLaman3}. For the converse, induct on
  the number of edges of $G$. If $G$ has a vertex $v$ of degree one
  with neighbor $u$, let $G' = G - E_G(v)$; then
  $G'^2 = G^2 - E_{G^2}(v)$ (\cref{lem:square-delete-degree-le-one}),
  $G'^2$ is Laman (\cref{lem:isLaman3-mono}), and $v$ has degree
  $d_G(u) \le 3$ in $G^2$ (\cref{lem:square-leaf-neighborhood},
  \cref{lem:isLaman3-degree-le-three}), so
  \cref{cor:zero-extension-degree-le-three} reduces the independence
  of $E(G^2)$ to that of $E(G'^2)$, which holds by induction.
  Otherwise every vertex has degree zero or at least two. If $G$ has
  no edges, $E(G^2)$ is empty. Else restrict to the support $S$ of
  $G$: squaring commutes with the restriction
  (\cref{lem:square-induce-support}), the Laman condition restricts
  (\cref{lem:isLaman3-induce}), and the restricted graph has minimum
  degree at least two on $S$, so \cref{thm:jacobs-min-degree-two}
  applies; \cref{lem:genericMatroid-induce-transport} transports the
  resulting independence back to $V$. This makes Jackson--Jord\'an's
  conditional Theorem~5.4 unconditional; the passage from maximum
  degree three to the minimum-degree-two core, asserted
  in~\cite{jacksonJordan2008} without proof, is the degree-one
  induction above.
\end{proof}

\subsection{The degree-one rank formula}
\label{sec:jacobs-degree-one}

The rank formula requires minimum degree at least two, and is false
without it: a single edge has $r(G^2) = 1$ while
$3|V| - 6 - \operatorname{def}(\tilde G) \le 0$. The correct statement
for graphs with vertices of degree one, due to
Franzblau~\cite{franzblau2000} for trees and Jackson--Jord\'an
(Lemma~4.2 of~\cite{jacksonJordan2008}) in general, reduces $r(G^2)$ to
the square of the \emph{core}, the maximal subgraph of minimum degree at
least two. Write $V_1(G)$ for the set of vertices of degree one in $G$.

\begin{definition}[The core]
  \label{def:two-core}
  The \emph{core} $G^{\mathrm{core}}$ of a graph $G$ is its maximal
  subgraph of minimum degree at least two: the join of all subgraphs
  $H \le G$ in which every vertex of the support of $H$ has
  $H$-degree at least two. It is empty if and only if every connected
  component of $G$ is a tree.
\end{definition}

\begin{theorem}[The degree-one rank formula for trees; JJ Lemma 4.2(a)]
  \label{thm:degree-one-rank-tree}
  \uses{def:genericRank, def:square-graph}
  Let $G$ be a tree on a finite vertex set $V$ with $|V| \ge 2$. Then
  \[
    r(G^2) \;=\; 2|V| - 5 + |V_1(G)| .
  \]
  This case is due to Franzblau~\cite{franzblau2000}.
\end{theorem}
\begin{proof}
  \uses{cor:zero-extension-clique-rank, lem:square-delete-degree-le-one,
    lem:genericMatroid-induce-transport, lem:square-leaf-neighborhood,
    lem:square-induce-support}
  Induction on $|V|$, with base case a single edge:
  $r(K_2^2) = r(K_2) = 1 = 2\cdot 2 - 5 + 2$. For $|V| \ge 3$, pick a
  leaf $v$ with neighbor $u$ and remove the edges at $v$: squaring
  commutes with the removal (\cref{lem:square-delete-degree-le-one}),
  and $v$ has $G^2$-degree $d_G(u)$ with square-neighborhood a clique
  (\cref{lem:square-leaf-neighborhood}), so the rank drops by
  $\min(3, d_G(u))$ (\cref{cor:zero-extension-clique-rank}). If
  $d_G(u) \ge 3$ the drop is $3$ and $|V_1|$ drops by one; if
  $d_G(u) = 2$ the drop is $2$ and $|V_1|$ is unchanged ($u$ becomes
  the new leaf). In both cases the right-hand side, evaluated for the
  tree $G - v$ on $|V| - 1$ vertices (\cref{lem:square-induce-support}
  and \cref{lem:genericMatroid-induce-transport} for the isolated
  vertex), drops by the same amount.
\end{proof}

\begin{theorem}[The degree-one rank formula; JJ Lemma 4.2(b)]
  \label{thm:degree-one-rank}
  \uses{def:two-core, def:genericRank, def:square-graph}
  Let $G$ be a connected simple graph on a finite vertex set $V$ with
  $|V| \ge 2$ that is not a tree. Then
  \[
    r(G^2) \;=\; r\bigl((G^{\mathrm{core}})^2\bigr)
      + 2\,\bigl|V \setminus V(G^{\mathrm{core}})\bigr|
      + |V_1(G)| ,
  \]
  where $V(G^{\mathrm{core}})$ is the support of the core.
\end{theorem}
\begin{proof}
  \uses{cor:zero-extension-clique-rank, lem:square-delete-degree-le-one,
    lem:square-leaf-neighborhood, lem:genericMatroid-induce-transport,
    lem:square-induce-support}
  Induction on $|V|$, as in \cref{thm:degree-one-rank-tree}: if
  $V_1(G) = \emptyset$ then $G = G^{\mathrm{core}}$ and both sides
  agree; otherwise remove the edges at a vertex $v$ of degree one with
  neighbor $u$ (here $d_G(u) \ge 2$, since $d_G(u) = 1$ would make the
  connected $G$ a single edge, a tree). The removed vertex has
  $G^2$-degree $d_G(u)$ with square-neighborhood a clique
  (\cref{lem:square-leaf-neighborhood}), so the rank drops by
  $\min(3, d_G(u))$ (\cref{cor:zero-extension-clique-rank};
  \cref{lem:square-delete-degree-le-one}). The core is unchanged in
  both cases; if $d_G(u) \ge 3$ the drop is $3$, matching a drop of
  one in $|V_1|$ and one non-core vertex fewer; if $d_G(u) = 2$ the
  drop is $2$, matching one non-core vertex fewer with $|V_1|$
  unchanged ($u$ becomes the new degree-one vertex). The induction
  closes at $V_1 = \emptyset$.
\end{proof}
